\section{Experimental Design}

The experiments use three tasks whose labels and group structure are fixed
before any policy is searched. CodeXGLUE defect detection provides 21{,}854
training, 2{,}732 dev and 2{,}732 test functions; the controller pools are drawn
only from the official training split, and functions are grouped by an exact
hash of their text so that identical functions cannot cross the internal
\texttt{search\_train}, \texttt{selection\_dev} and \texttt{confirm} pools
\citep{lu2021codexglue,zhou2019devign}. Each internal pool contains 60
functions. This controls duplicate-function leakage. FinDVer provides financial claims with \texttt{supported}/\texttt{not
supported} labels; its pools are grouped by report, and contain 60, 59 and 44
claims \citep{zhao2024findver}. The official FinDVer test labels are not used
for controller search. ProcessBench contains 3{,}400 annotated records over
2{,}842 mathematical problems; its pools are grouped by problem and contain 60
records each \citep{zheng2025processbench,cobbe2021gsm8k,lightman2024verify,uesato2022process}.
FinDVer internal pools are carved from official testmini and grouped by report
(41/16/16 groups; 60/59/44 claims); labels are used inside these pools for
fitting, thresholding and frozen evaluation. The official test split is
separate and read only by the locked-test evaluator. State is a \texttt{Claim:}
line, a blank line, then an \texttt{Evidence:} block. CodeXGLUE state is raw
function text
(positive rates 0.50/0.40/0.45). ProcessBench asks ``Which step contains the
earliest error?''; its state is the problem plus \texttt{[k] step} lines and an
\texttt{ALL STEPS CORRECT} sentinel.

All inference runs on one Windows machine with CPU-only PyTorch in FP32, two
threads and one model resident at a time. Laya is pinned to source revision
\texttt{b9fe351} and checkpoint revision \texttt{7b928d8}; Vela is pinned to
revision \texttt{d6f03aa}. Peak resident memory stays below 2.2~GB. The native
single-action baselines on \texttt{search\_dev} are as follows: CodeXGLUE
accuracy is 0.450/0.650 for Laya
(\noul/\choice) and 0.633/0.667 for Vela; FinDVer is 0.712/0.644 for Laya and
0.458/0.458 for Vela; ProcessBench \choice{} accuracy is 0.10 for both models.
The near-identical per-call latency of the two representations (about
2.2--2.8~s) is why call count, rather than a model-specific latency model, is
the cost term in Equation~\ref{eq:fitness}. The native numbers already contain a warning: the better representation depends
on the task and the model, and on FinDVer the two Vela actions are
indistinguishable at the aggregate level.

The first experiment measures action complementarity without training a router. For each
task--model cell we take the better fixed action on \texttt{search\_train} and
freeze that direction; we then count rows where the alternative action repairs a
base error, rows where it breaks a correct base answer, and their difference.
The oracle headroom is the accuracy of a per-row oracle minus the accuracy of
the better fixed action. Because the direction is fixed on train, the held-out
net-repair numbers are not selected on the evaluation pool. We report net repair
in percentage points with a group-level percentile bootstrap (5{,}000
resamples), and pair it with the macro-F1 difference between fixed \choice{} and
fixed \noul{} on the same rows. ProcessBench has only one evaluated
representation, so its headroom is zero by construction and it is excluded from
the binary complementarity table.

The second experiment asks whether the signed incremental value of the second
action is predictable before that action is executed. The target is
$\Delta_i=\mathbf{1}[\choice_i\text{ correct}]-
\mathbf{1}[\noul_i\text{ correct}]\in\{-1,0,+1\}$. The features are deployable at the
moment of decision: $p$, $|p-0.5|$ and the binary entropy $-p\log p -
(1-p)\log(1-p)$; no \choice{} output, gold label or downstream feature is used.
We compare logistic regression, a shallow decision tree, a single confidence
threshold and a never-switch majority baseline. Training uses group-aware
three-fold out-of-fold predictions; the predictor and threshold are then chosen
together on \texttt{selection\_dev} by cost-adjusted net utility, frozen, and
evaluated once on \texttt{confirm}. Net utility is (repairs $-$ regressions)
divided by the number of rows; the cost-adjusted variant subtracts
$\lambda=0.05$ per switch. A champion that switches every row is retained, because it reveals
whether the predictor adds anything beyond fixed switching.

The search comparison runs optimisers under a matched evaluation budget. GA, ES and NSGA-II
are run on the cached pools with 24 unique genome evaluations and seeds 17, 23
and 41; the five-seed extension adds 59 and 73. Random search receives the same
budget. The primary deployment comparison is pre-specified: the GA genome
selected on \texttt{search\_train} is frozen and compared with fixed
switch-only on \texttt{confirm} using paired group-bootstrap macro-F1 intervals.
All other cell-level comparisons are exploratory and are reported with effect
estimates and intervals rather than as significance claims. We then examine the
search trajectory itself: best-so-far fitness per generation, the number of
unique genomes actually evaluated, the fraction of single-gene mutations that
leave fitness unchanged, the rank correlation between train and dev fitness
across the archive, and the effect of raising the budget cap from 24 to 48 on
CodeXGLUE with Vela.

Two boundary checks test whether the conclusions generalise. The transfer experiment freezes each source-task
genome and applies it to every target task and pool without re-searching;
diagonal cells are same-task deployment checks, off-diagonal cells are true
transfer. The fault map partitions rows by task, model and a confidence band
derived from \noul{} probability or candidate margin, and reports the error
severity of each cell on discover, dev and confirm. A cell is eligible when it
contains at least eight rows on the discover pool, and it is confirmed when the
dev or confirm severity is within 0.15 of the discover severity.

Two statistical details matter for reading the tables. All binary comparisons
use paired group bootstrap with 5{,}000 resamples and a fixed random state
\citep{efron1993bootstrap}; the resampling unit is the group (function hash or
report), not the row, so dependence within a report is preserved. We report
McNemar exact tests for repair/harmful counts \citep{mcnemar1947} and Wilson
intervals for within-cell severity \citep{wilson1927}, but we do not call a
result significant on the basis of a single seed or a single pool. The locked
test was run once under a protocol frozen before any test call. The confirm
pools are held out from search and threshold selection but have been inspected
during analysis.
